\documentclass[11pt,a4paper]{article}

% Packages
\usepackage[utf8]{inputenc}
\usepackage[T1]{fontenc}
\usepackage{booktabs}
\usepackage{graphicx}
\usepackage{subcaption}
\usepackage[margin=1in]{geometry}
\usepackage{hyperref}
\usepackage{longtable}

% Title information
\title{ORTPS Analysis}
\author{Data, Policy and Innovation Centre}
\date{\today}

\begin{document}

\maketitle
\tableofcontents
\newpage

\listoftables
\listoffigures
\newpage

\section{Data \& Methodology}

\subsection{Research Objective}

This analysis investigates the representation of Odisha Right to Public Services (ORTPS) in the state's grievance redressal system. We address two key questions:
\begin{enumerate}
    \item Do ORTPS-covered services appear as distinct categories in citizen complaints submitted through the grievance redressal portal?
    \item If so, what issues do complaints related to these services typically address?
\end{enumerate}

Understanding ORTPS representation in grievance data can reveal service delivery bottlenecks, inform resource allocation decisions, and identify priority areas for process improvements.

\subsection{Methodological Challenge}

The primary challenge in identifying ORTPS-related complaints lies in the mismatch between administrative categorization and service-based classification. The existing complaint category system was not designed with the ORTPS framework in mind---categories reflect organizational structure (departments, divisions) rather than specific public services delivered to citizens.

For example, a complaint about a delayed caste certificate might be categorized under the Education category administratively, since a caste certificate may be needed to avail reserved seats in colleges and universities, but from a service delivery perspective, it falls under the ORTPS framework's identity documentation services. This misalignment necessitates an alternative approach to identify ORTPS-related complaints based on the complaint text itself rather than pre-assigned categories.

\subsection{ORTPS Category Consolidation}

We began with the full list of ORTPS-covered services as provided by CAGR and consolidated them into 8 issue-based categories covering major service domains. This consolidation balances specificity with analytical tractability, grouping related services that citizens and administrators perceive as similar types of requests.

The 8 consolidated categories are:

\begin{enumerate}
    \item \textbf{Ration card} --- Issuance of new ration cards, inclusion/exclusion of household members, card corrections, and related PDS card-service requests (AAY/BPL/PHH/Food Security)

    \item \textbf{Certificates} --- Caste/resident/income/EWS certificates, birth and death certificates, marriage registration, legal heir certificates, and related certificate correction/reissue requests

    \item \textbf{Land matters} --- RoR and mutation issues, encumbrance/registration document access, land conversion and partition requests, and related land administration services

    \item \textbf{Scholarship} --- Pre-matric, post-matric, and state scholarship application, verification, sanction, and payment-related grievances

    \item \textbf{Utilities \& Connections} --- Piped water and non-industrial electricity connection requests and related service activation issues

    \item \textbf{Police \& Legal} --- Verification/legal service requests (including police/legal documentation and attestation-related services) and associated grievance follow-up

    \item \textbf{Vehicle Services} --- Driving licences (new/renewal/learner/address change), vehicle registration certificate (RC) copies, ownership transfer, and related transport-document services

    \item \textbf{Building \& Construction} --- Building plan approvals by ULBs/development authorities, building alteration permissions, and related approval workflows
\end{enumerate}

\subsection{Text-Based Classification Pipeline}

Given the administrative categorization challenge, we developed a two-stage automated classification pipeline operating directly on complaint text.

\subsubsection{Language Detection}

To ensure reliable keyword matching, we first filter complaints to English-language text using a two-stage detection approach:

\begin{enumerate}
    \item \textbf{Script Detection:} A regex-based filter identifies and excludes complaints containing Odia or Devanagari (Hindi) scripts, which are detectable by their Unicode character ranges.

    \item \textbf{Lingua Confidence Scoring:} For complaints passing the script filter, we use the Lingua language detection library to compute English-language confidence scores. Complaints with confidence $\geq 0.85$ are classified as English and retained for further analysis.
\end{enumerate}

This two-stage approach balances precision and recall, filtering out non-English text while preserving code-mixed complaints where English is dominant.

\subsubsection{Category Labeling}

For English-language complaints, we employ a hybrid classification approach combining keyword matching and semantic similarity:

\begin{enumerate}
    \item \textbf{Keyword Matching:} Case-insensitive substring matching against category-specific keyword lists. For example, ``Scholarship'' includes keywords such as \textit{scholarship}, \textit{post matric scholarship}, \textit{pre matric}, and related terms. This method captures complaints with explicit mentions of service names.

    \item \textbf{Embedding Similarity:} Complaints not matched by keywords are evaluated using sentence embeddings (all-MiniLM-L6-v2 model) to compute semantic similarity to category prototypes. Complaints with cosine similarity $\geq 0.60$ to any category are assigned to that category. This method captures semantically related complaints that may not use exact service terminology (e.g., ``my post matric payment is pending'' $\rightarrow$ Scholarship).
\end{enumerate}

Complaints matching neither method remain unlabeled, indicating they likely concern non-ORTPS services (e.g., infrastructure, law and order, general administration).

\subsection{Data Scope}

We restrict our analysis along two dimensions:

\begin{itemize}
    \item \textbf{Temporal scope:} July 2023 onwards (fiscal years 2023--2025). This recent time window reflects current service delivery patterns and ensures sufficient data volume for category-level analysis.

    \item \textbf{Language filter:} English-only complaints. Since our ORTPS keyword lists were developed in English, language filtering ensures classification reliability. While this excludes Odia-language complaints, English remains the predominant language in the portal (over 70\% of total complaints).
\end{itemize}

The final analysis dataset comprises \textbf{763,237 English-language complaints} from July 2023 onwards.

\section{Results: ORTPS Representation}

\begin{table}[htbp]
\caption{ORTPS Category Aggregation by Fiscal Year}
\label{tab:ortps_aggregation}
\resizebox{\textwidth}{!}{%
\begin{tabular}{lrrrr}
\toprule
Fiscal Year & \multicolumn{2}{c}{2023} & \multicolumn{2}{c}{2024} \\
\cmidrule(lr){2-3} \cmidrule(lr){4-5}
Metric & Count & % of Total & Count & % of Total \\
Category &  &  &  &  \\
\midrule
Ration card & 8,863 & 3.25\% & 14,505 & 3.22\% \\
Certificates & 2,234 & 0.82\% & 3,829 & 0.85\% \\
Land matters & 540 & 0.20\% & 2,066 & 0.46\% \\
Scholarship & 1,025 & 0.38\% & 1,744 & 0.39\% \\
Utilities \& Connections & 625 & 0.23\% & 1,187 & 0.26\% \\
Police \& Legal & 161 & 0.06\% & 311 & 0.07\% \\
Vehicle Services & 135 & 0.05\% & 224 & 0.05\% \\
Building \& Construction & 55 & 0.02\% & 92 & 0.02\% \\
\bottomrule
\end{tabular}
}
\end{table}




\subsection{Temporal Trends and Cross-Sectional Variation}

Table~\ref{tab:ortps_aggregation} presents the fiscal year aggregation of ORTPS-related grievances, showing both temporal trends and cross-sectional variation across our pre-defined categories.

Across the 8 updated categories, total ORTPS-identified complaints increased from \textbf{13,638 in FY2023} to \textbf{23,958 in FY2024} (\textbf{+75.7\%}).

Key cross-sectional patterns are:

\begin{itemize}
    \item \textbf{Ration card} remains the largest category by a wide margin (8,863 $\rightarrow$ 14,505 complaints), with a stable share of total grievances (3.25\% in FY2023 vs 3.22\% in FY2024). This indicates persistent high demand for ration-card service delivery.

    \item \textbf{Certificates} is the second-largest category (2,234 $\rightarrow$ 3,829; +71.4\%), with share increasing from 0.82\% to 0.85\%. Certificate-linked grievances are both high-volume and growing.

    \item \textbf{Land matters} shows the fastest relative growth (540 $\rightarrow$ 2,066; nearly 4x), with share rising from 0.20\% to 0.46\%. This is the sharpest year-on-year expansion among the major categories.

    \item \textbf{Scholarship} also grows materially in absolute terms (1,025 $\rightarrow$ 1,744; +70.1\%) while share remains stable (0.38\% to 0.39\%), suggesting sustained demand that scales with overall complaint growth.

    \item \textbf{Smaller categories} (Utilities \& Connections, Police \& Legal, Vehicle Services, Building \& Construction) all increase in count between FY2023 and FY2024, but each remains below 0.30\% share, indicating long-tail service demand relative to the four dominant categories above.
\end{itemize}

\subsection{Word Clouds for Four Focus Categories}

The figures below show high-frequency terms for FY2023-24 and FY2024-25 in the four categories analyzed in the following section.

\begin{figure}[htbp]
\centering
\begin{subfigure}[b]{0.48\textwidth}
    \centering
    \includegraphics[width=\textwidth]{wordclouds/Ration_Card_FY2023_2024.png}
    \caption{FY 2023-24}
    \label{fig:ration_card_focus_2023}
\end{subfigure}
\hfill
\begin{subfigure}[b]{0.48\textwidth}
    \centering
    \includegraphics[width=\textwidth]{wordclouds/Ration_Card_FY2024_2025.png}
    \caption{FY 2024-25}
    \label{fig:ration_card_focus_2024}
\end{subfigure}
\caption{Ration card Word Clouds}
\end{figure}

\begin{figure}[htbp]
\centering
\begin{subfigure}[b]{0.48\textwidth}
    \centering
    \includegraphics[width=\textwidth]{wordclouds/Certificates_FY2023_2024.png}
    \caption{FY 2023-24}
    \label{fig:certificates_focus_2023}
\end{subfigure}
\hfill
\begin{subfigure}[b]{0.48\textwidth}
    \centering
    \includegraphics[width=\textwidth]{wordclouds/Certificates_FY2024_2025.png}
    \caption{FY 2024-25}
    \label{fig:certificates_focus_2024}
\end{subfigure}
\caption{Certificates Word Clouds}
\end{figure}

\begin{figure}[htbp]
\centering
\begin{subfigure}[b]{0.48\textwidth}
    \centering
    \includegraphics[width=\textwidth]{wordclouds/Land_matters_FY2023_2024.png}
    \caption{FY 2023-24}
    \label{fig:land_matters_focus_2023}
\end{subfigure}
\hfill
\begin{subfigure}[b]{0.48\textwidth}
    \centering
    \includegraphics[width=\textwidth]{wordclouds/Land_matters_FY2024_2025.png}
    \caption{FY 2024-25}
    \label{fig:land_matters_focus_2024}
\end{subfigure}
\caption{Land matters Word Clouds}
\end{figure}

\begin{figure}[htbp]
\centering
\begin{subfigure}[b]{0.48\textwidth}
    \centering
    \includegraphics[width=\textwidth]{wordclouds/Scholarship_FY2023_2024.png}
    \caption{FY 2023-24}
    \label{fig:scholarship_focus_2023}
\end{subfigure}
\hfill
\begin{subfigure}[b]{0.48\textwidth}
    \centering
    \includegraphics[width=\textwidth]{wordclouds/Scholarship_FY2024_2025.png}
    \caption{FY 2024-25}
    \label{fig:scholarship_focus_2024}
\end{subfigure}
\caption{Scholarship Word Clouds}
\end{figure}

\section{Results: ORTPS Nature of complaints}

Using the updated category scheme in Table~\ref{tab:ortps_aggregation}, we reviewed topic-model outputs and labeled topic samples for four priority categories: \textbf{Ration card}, \textbf{Certificates}, \textbf{Land matters}, and \textbf{Scholarship}. The goal here is descriptive: what these ORTPS-identified complaints actually contain.

\subsection{Category-Specific Complaint Content}

\subsubsection{Ration card (FY2023: 8,863; FY2024: 14,505)}

Topic modeling on 11,910 complaints (29 topics; 72.5\% modeled coverage) shows that the largest single topic is \textbf{PM Awas housing requests} (Topic 0: 1,101 complaints; 9.2\%). The first clearly ration-card-specific cluster appears at Topic 12 (272 complaints; 2.3\%), with keywords centered on ``ration card'', ``new ration card'', and ``add''.

Labeled samples from the ration-card topic indicate genuine ORTPS service problems:
\begin{itemize}
    \item long-pending card issuance despite prior applications;
    \item failure to add family members to existing cards;
    \item online record-sync and Aadhaar-duplicate errors blocking service access;
    \item repeated non-approval at block/MI level.
\end{itemize}

\textbf{Assessment:} The category contains authentic ration-card delivery bottlenecks, but the overall pool is heavily diluted by non-ORTPS housing petitions, generic escalation templates, and unrelated sectoral grievances.

\subsubsection{Certificates (FY2023: 2,234; FY2024: 3,829)}

Topic modeling on 3,192 complaints (18 topics; 85.2\% modeled coverage) again shows \textbf{PM Awas} as the largest topic (Topic 0: 620 complaints; 19.4\%). The explicit certificate-services topic is Topic 5 (144 complaints; 4.5\%), with keywords including ``caste certificate'', ``birth'', and ``legal heir certificate''.

Labeled certificate samples show real ORTPS certificate issues:
\begin{itemize}
    \item caste/resident/income certificate processing delays beyond due date;
    \item legal-heir certificate approval/rejection errors;
    \item birth/DOB correction requests affecting exam records;
    \item rejection logic disputes (e.g., ROR-related objections for landless applicants).
\end{itemize}

\textbf{Assessment:} A real certificate-service signal is present, but it is not dominant. The category operates as a mixed bucket with substantial housing, governance, and generic grievance noise.

\subsubsection{Land matters (FY2023: 540; FY2024: 2,066)}

Topic modeling on 244 complaints (2 topics; 100.0\% modeled coverage) is highly concentrated: Topic 0 contains 93 complaints (38.1\%) and combines ``house/yojana/pradhan/mantri'' language; Topic 1 (65 complaints; 26.6\%) captures ``plot/ror/record'' patterns alongside non-land administrative noise.

Labeled samples show mixed content:
\begin{itemize}
    \item legitimate land-service issues (mutation status, ROR update mismatches, encumbrance certificate portal failures);
    \item non-land content (PMAY housing requests, ration-card member addition, police/FIR and corruption complaints).
\end{itemize}

\textbf{Assessment:} The current Land matters label has weak thematic purity. It includes some high-value land administration complaints, but category-level counts are materially inflated by off-topic requests.

\subsubsection{Scholarship (FY2023: 1,025; FY2024: 1,744)}

Topic modeling on 1,475 complaints (9 topics; 81.3\% modeled coverage) shows PM Awas as the largest cluster (Topic 0: 375 complaints; 25.4\%). The scholarship-specific topic is Topic 1 (183 complaints; 12.4\%), with terms centered on scholarship application, approval, and payment.

Labeled scholarship-topic samples capture genuine scholarship service requests:
\begin{itemize}
    \item non-sanction or delayed sanction after institutional/district verification;
    \item missing instalments after partial disbursement;
    \item eligibility and scheme-inclusion disputes (including Nua-O and Gopabandhu scholarship contexts).
\end{itemize}

\textbf{Assessment:} True scholarship grievances are clearly observable, but they are not the dominant theme in the category-level corpus.

\subsection{Cross-Cutting Pattern in the Four Focus Categories}

Across all four categories, the same pattern recurs: \textbf{PM Awas housing requests appear as Topic 0 in every category}. This indicates systematic lexical overlap between ORTPS service keywords and broader welfare/infrastructure language (``house'', ``awas'', ``construction''), which drives false positives.

In addition, a second recurrent noise source is template-driven complaints (mass corruption allegations and ``take necessary action'' forwarding text), which further reduces category specificity.

\subsection{Implication for Interpreting ORTPS Counts}

For these four categories, reported volumes in Table~\ref{tab:ortps_aggregation} should be treated as \emph{upper-bound estimates} of true service-specific grievance demand unless the pipeline explicitly filters:
\begin{enumerate}
    \item non-ORTPS housing scheme terms (PMAY/Biju Pakka Ghar variants), and
    \item template-style escalation/corruption text not tied to specific ORTPS transactions.
\end{enumerate}

\newpage

%\section{Main Exhibits}

% \textbf{Key Findings (FY2023-2025):}
% \begin{itemize}
%     \item \textbf{Ration Card dominates ORTPS grievances:} Consistently the highest volume category across all fiscal years, with 8,593 complaints (3.13\% of total) in FY2023 growing to 15,052 (3.28\%) in FY2024---a 75\% year-over-year increase. This category alone represents over 3\% of all ORTPS-related grievances.

%     \item \textbf{Income Certificate shows dramatic FY2024 surge:} Volume increased 4.6$\times$ from FY2023 to FY2024 (338 $\rightarrow$ 1,542 complaints), with share nearly tripling from 0.12\% to 0.34\%. This sharp discontinuity warrants investigation into potential systemic issues, policy changes, or portal disruptions that emerged in FY2024.

%     \item \textbf{Scholarship complaints show sustained growth:} Volume increased 76\% from FY2023 to FY2024 (1,074 $\rightarrow$ 1,895), with share rising from 0.39\% to 0.41\%. This trend suggests increasing demand for scholarship services and growing awareness among beneficiaries.

%     \item \textbf{Caste Certificate growth lags overall trends:} While absolute numbers increased 27\% from FY2023 to FY2024 (600 $\rightarrow$ 760), the share declined from 0.22\% to 0.17\%, indicating this service is growing slower than overall ORTPS complaint volume.

%     \item \textbf{FY2025 data is partial (incomplete):} All categories show lower absolute counts relative to FY2023 and FY2024, but some percentage shares are elevated (e.g., Scholarship at 0.79\% vs. 0.41\% in FY2024), indicating that data collection for FY2025 is incomplete as the fiscal year is ongoing.
% \end{itemize}

\newpage

\section{Exploratory Analysis}

\subsection{Fiscal Year Aggregation}

This table presents the aggregation of ORTPS-related grievances by fiscal year (July-June) and category for FY2023 onwards.

\begin{table}[htbp]
\caption{ORTPS Category Aggregation by Fiscal Year}
\label{tab:ortps_aggregation}
\resizebox{\textwidth}{!}{%
\begin{tabular}{lrrrr}
\toprule
Fiscal Year & \multicolumn{2}{c}{2023} & \multicolumn{2}{c}{2024} \\
\cmidrule(lr){2-3} \cmidrule(lr){4-5}
Metric & Count & % of Total & Count & % of Total \\
Category &  &  &  &  \\
\midrule
Ration card & 8,863 & 3.25\% & 14,505 & 3.22\% \\
Certificates & 2,234 & 0.82\% & 3,829 & 0.85\% \\
Land matters & 540 & 0.20\% & 2,066 & 0.46\% \\
Scholarship & 1,025 & 0.38\% & 1,744 & 0.39\% \\
Utilities \& Connections & 625 & 0.23\% & 1,187 & 0.26\% \\
Police \& Legal & 161 & 0.06\% & 311 & 0.07\% \\
Vehicle Services & 135 & 0.05\% & 224 & 0.05\% \\
Building \& Construction & 55 & 0.02\% & 92 & 0.02\% \\
\bottomrule
\end{tabular}
}
\end{table}


\subsection{Word Cloud Analysis}

The following word clouds visualize the most frequent terms in grievance text for each ORTPS category for FY2023-24 and FY2024-25.

% Caste Certificate
\begin{figure}[htbp]
\centering
\begin{subfigure}[b]{0.48\textwidth}
    \centering
    \includegraphics[width=\textwidth]{wordclouds/Caste_certificate_FY2023_2024.png}
    \caption{FY 2023-24}
    \label{fig:caste_certificate_2023}
\end{subfigure}
\hfill
\begin{subfigure}[b]{0.48\textwidth}
    \centering
    \includegraphics[width=\textwidth]{wordclouds/Caste_certificate_FY2024_2025.png}
    \caption{FY 2024-25}
    \label{fig:caste_certificate_2024}
\end{subfigure}

\caption{Caste Certificate Grievances}
\label{fig:caste_certificate}
\end{figure}

\textbf{Note:} Word frequency analysis for FY2023-25 reveals common pain points in Caste Certificate processing. Look for terms related to application status, delays, document requirements, and verification processes. Cross-year comparison can identify persistent bottlenecks versus emerging issues.

% Income Certificate
\begin{figure}[htbp]
\centering
\begin{subfigure}[b]{0.48\textwidth}
    \centering
    \includegraphics[width=\textwidth]{wordclouds/Income_certificate_FY2023_2024.png}
    \caption{FY 2023-24}
    \label{fig:income_certificate_2023}
\end{subfigure}
\hfill
\begin{subfigure}[b]{0.48\textwidth}
    \centering
    \includegraphics[width=\textwidth]{wordclouds/Income_certificate_FY2024_2025.png}
    \caption{FY 2024-25}
    \label{fig:income_certificate_2024}
\end{subfigure}

\caption{Income Certificate Grievances}
\label{fig:income_certificate}
\end{figure}

\textbf{Note:} Given the 4.6$\times$ spike in FY2024 (Table \ref{tab:ortps_aggregation}), compare word clouds across years to identify new terms or increased frequency of specific issues that may explain this surge. Pay attention to portal-related terminology or policy changes.

% Scholarship
\begin{figure}[htbp]
\centering
\begin{subfigure}[b]{0.48\textwidth}
    \centering
    \includegraphics[width=\textwidth]{wordclouds/Scholarship_FY2023_2024.png}
    \caption{FY 2023-24}
    \label{fig:scholarship_2023}
\end{subfigure}
\hfill
\begin{subfigure}[b]{0.48\textwidth}
    \centering
    \includegraphics[width=\textwidth]{wordclouds/Scholarship_FY2024_2025.png}
    \caption{FY 2024-25}
    \label{fig:scholarship_2024}
\end{subfigure}

\caption{Scholarship Grievances}
\label{fig:scholarship}
\end{figure}

\textbf{Note:} Scholarship complaints show consistent 70.1\% growth from FY2023-24. Word clouds may reveal seasonal patterns (academic calendar dependencies), common delays in disbursement, or eligibility verification issues. Terms related to deadlines and urgency are likely prominent.

% Ration Card
\begin{figure}[htbp]
\centering
\begin{subfigure}[b]{0.48\textwidth}
    \centering
    \includegraphics[width=\textwidth]{wordclouds/Ration_Card_FY2023_2024.png}
    \caption{FY 2023-24}
    \label{fig:ration_card_2023}
\end{subfigure}
\hfill
\begin{subfigure}[b]{0.48\textwidth}
    \centering
    \includegraphics[width=\textwidth]{wordclouds/Ration_Card_FY2024_2025.png}
    \caption{FY 2024-25}
    \label{fig:ration_card_2024}
\end{subfigure}

\caption{Ration Card Grievances}
\label{fig:ration_card}
\end{figure}

\textbf{Note:} As the dominant ORTPS category with 63.7\% FY2023-24 growth, Ration Card word clouds are critical for understanding systemic service delivery challenges. Expect high frequency of terms related to updates, corrections, additions/deletions of members, and digitization issues (e.g., Aadhaar linking).

\end{document}
